\documentclass[12pt, a4paper]{article}
\usepackage[utf8]{inputenc}
\usepackage[spanish]{babel}
\usepackage{geometry}
\usepackage{array}
\usepackage[normalem]{ulem} % Paquete para subrayados dobles/triples

% Configuración de la página
\geometry{
    a4paper,
    left=2.5cm,
    right=2.5cm,
    top=2.5cm,
    bottom=2.5cm
}

% --- Comandos para formato de claves ---
\newcommand{\pk}[1]{\underline{\texttt{#1}}}      % Para Primary Key (1 subrayado)
\newcommand{\fk}[1]{\uuline{\texttt{#1}}}      % Para Foreign Key (2 subrayados)
\newcommand{\pkc}[1]{\uuline{\underline{\texttt{#1}}}} % Para Primary Key Compuesta (3 subrayados)


\begin{document}

% --- PORTADA ---
\begin{titlepage}
    \centering
    \vspace*{\stretch{1}}
    {\Huge\bfseries Normalización de Base de Datos}
    \vspace{1.5cm}
    \vfill
    {\Large\bfseries Asignatura:} \\
    {\large Base de Datos Estructuradas}
    \vspace{2cm}
    \vfill
    {\Large\bfseries Autor:} \\
    {\large Andres Lopez Gonzalez}
    \vspace*{\stretch{2}}\\
    {\large \today}
\end{titlepage}

\newpage

% --- INICIO DEL DOCUMENTO ---

\section*{Tabla Inicial sin Normalizar}

Se parte de la siguiente tabla que contiene información sobre productos, sus materias primas y los proveedores (es una sola fila, pero por temas de espacio en la hoja se muestrea como dos).

\begin{center}
\begin{tabular}{|c|c|c|c|}
\hline
\texttt{codprod} & \texttt{codprove} & \texttt{codmate} & \texttt{nomprov} \\
\hline
 & & & \\
\hline
\end{tabular}
\vspace{0.1cm} % Espacio pequeño entre las tablas
\begin{tabular}{|c|c|c|c|}
\hline
\texttt{dirprov} & \texttt{descprod} & \texttt{descmate} & \texttt{valorprod} \\
\hline
 & & & \\
\hline
\end{tabular}
\end{center}

La estructura inicial, en una sola tabla, es la siguiente:
\begin{center}
\textbf{PRODUCTO\_GENERAL}(\pk{codprod}, \texttt{descprod}, \texttt{valorprod}, \fk{codprove}, \texttt{nomprov}, \texttt{dirprov}, \fk{codmate}, \texttt{descmate})
\end{center}

El objetivo es normalizar esta estructura hasta la Tercera Forma Normal (3NF) para eliminar redundancias y anomalías de datos.

\section{Primera Forma Normal (1NF)}
\textit{\textbf{Definición:} Una tabla está en 1NF si todos sus atributos son atómicos, es decir, cada celda contiene un único valor y no existen grupos de repetición.}

La tabla inicial ya cumple con la Primera Forma Normal, ya que cada columna representa un atributo único y no hay campos que almacenen múltiples valores, dicho de otra forma, la función que establece esa 'casilla' no es multievaluada.

La estructura en 1NF es la misma que la tabla original, identificando la clave primaria:
\begin{center}
\textbf{PRODUCTO\_1NF}(\pk{codprod}, \texttt{descprod}, \texttt{valorprod}, \fk{codprove}, \texttt{nomprov}, \texttt{dirprov}, \fk{codmate}, \texttt{descmate})
\end{center}

\section{Segunda Forma Normal (2NF)}
\textit{\textbf{Definición:} Una tabla está en 2NF si está en 1NF y todos los atributos no clave dependen completamente de la clave primaria completa.}

Esta regla es relevante principalmente para tablas con claves primarias compuestas (formadas por más de una columna). En nuestro caso, la clave primaria es \pk{codprod}, que es una clave simple.

Dado que la clave primaria no es compuesta, no pueden existir dependencias parciales. Por lo tanto, la tabla en 1NF satisface automáticamente la Segunda Forma Normal.

La estructura en 2NF es idéntica a la de 1NF:
\begin{center}
\textbf{PRODUCTO\_2NF}(\pk{codprod}, \texttt{descprod}, \texttt{valorprod}, \fk{codprove}, \texttt{nomprov}, \texttt{dirprov}, \fk{codmate}, \texttt{descmate})
\end{center}

\section{Tercera Forma Normal (3NF)}
\textit{\textbf{Definición:} Una tabla está en 3NF si está en 2NF y no contiene dependencias transitivas. Una dependencia transitiva ocurre cuando un atributo no clave depende de otro atributo no clave.}

Al analizar la tabla en 2NF, identificamos las siguientes dependencias transitivas:

\begin{itemize}
    \item \texttt{nomprov} (Nombre del Proveedor) y \texttt{dirprov} (Dirección del Proveedor) dependen de \texttt{codprove} (Código de Proveedor), que no es la clave primaria.
    
    \textbf{Dependencia:} \pk{codprod} $\rightarrow$ \fk{codprove} $\rightarrow$ (\texttt{nomprov}, \texttt{dirprov})

    \item \texttt{descmate} (Descripción de Materia Prima) depende de \texttt{codmate} (Código de Materia Prima), que no es la clave primaria.
    
    \textbf{Dependencia:} \pk{codprod} $\rightarrow$ \fk{codmate} $\rightarrow$ \texttt{descmate}
\end{itemize}

Para eliminar estas dependencias transitivas, dividimos la tabla en tres tablas separadas, una para cada entidad: Productos, Proveedores y Materias Primas; la estructura final en Tercera Forma Normal consta de las siguientes tres tablas:

\vspace{0.5cm}
\begin{center}
    \textbf{1. Tabla de Productos}
    
    \vspace{0.2cm}
    \begin{tabular}{|l|l|}
    \hline
    \multicolumn{2}{|c|}{\textbf{PRODUCTOS}} \\
    \hline
    \pk{codprod} & Clave Primaria \\
    \texttt{descprod} & \\
    \texttt{valorprod} & \\
    \fk{codprove} & Clave Foránea a PROVEEDORES \\
    \fk{codmate} & Clave Foránea a MATERIAS\_PRIMAS \\
    \hline
    \end{tabular}
\end{center}

\vspace{0.5cm}
\begin{center}
    \textbf{2. Tabla de Proveedores}

    \vspace{0.2cm}
    \begin{tabular}{|l|l|}
    \hline
    \multicolumn{2}{|c|}{\textbf{PROVEEDORES}} \\
    \hline
    \pk{codprove} & Clave Primaria \\
    \texttt{nomprov} & \\
    \texttt{dirprov} & \\
    \hline
    \end{tabular}
\end{center}

\vspace{0.5cm}
\begin{center}
    \textbf{3. Tabla de Materias Primas}

    \vspace{0.2cm}
    \begin{tabular}{|l|l|}
    \hline
    \multicolumn{2}{|c|}{\textbf{MATERIAS\_PRIMAS}} \\
    \hline
    \pk{codmate} & Clave Primaria \\
    \texttt{descmate} & \\
    \hline
    \end{tabular}
\end{center}
\vspace{1cm}

Con esta estructura, se eliminan las redundancias de datos. La información de cada proveedor y de cada materia prima se almacena una sola vez, y las tablas se relacionan mediante claves foráneas. El sistema está ahora en 3NF.

\end{document}